\begin{ouverture}[geometrie]{16}{Polygones, du triangle à l’hexagone}{Partie 4 · Espace et géométrie}
{Triangles · Quadrilatères · Trapèzes · Pentagones · Hexagones}
\objectif{Reconnaître et nommer les polygones jusqu’à l’hexagone.}
\objectif{Utiliser les propriétés des figures usuelles, trapèzes compris.}
\objectif{Construire ces figures et justifier avec les instruments.}
\prerequis{Constructions du chapitre 15.}
\end{ouverture}
\pagestyle{otMain}\essentiel\fontsize{11}{15}\selectfont
\phantomsection\label{cours-16-1}
\essSection{1}{Les triangles, bien en main}
\begin{definition}Un polygone est une figure fermée dont les côtés sont des segments. Trois côtés : triangle ; quatre : quadrilatère ; cinq : pentagone ; six : hexagone. Le triangle rectangle a un angle droit, l’isocèle deux côtés égaux, l’équilatéral trois.\end{definition}
\label{fig-cours-16-1-cours-triangles}\figureGeo{cours-triangles}

\begin{propriete}Dans un triangle isocèle, les angles à la base sont égaux ; dans un équilatéral, les trois angles sont égaux. Les codes affirment ces propriétés sur la figure.\end{propriete}
\begin{savoirfaire}[Le compas reporte les côtés]Construire au compas : pour un isocèle de base 4 cm et de côtés 6 cm, tracer la base, puis deux arcs de rayon 6 cm centrés aux extrémités ; leur croisement donne le sommet principal.\end{savoirfaire}
\clearpage
\phantomsection\label{cours-16-2}
\essSection{2}{Quadrilatères et trapèzes}
\begin{definition}Le rectangle a quatre angles droits ; le losange, quatre côtés égaux ; le carré, les deux. Le trapèze a deux côtés opposés parallèles ; le trapèze rectangle possède en plus deux angles droits.\end{definition}
\label{fig-cours-16-2-cours-trapeze}\figureGeo{cours-trapeze}

\begin{propriete}Dans le rectangle, le losange et le carré, les côtés opposés sont parallèles deux à deux. Dans le trapèze, seule la paire de côtés parallèles est garantie : les deux autres côtés peuvent être quelconques.\end{propriete}
\begin{attention}Un carré est à la fois un rectangle et un losange : il en possède toutes les propriétés. Une figure peut porter plusieurs noms ; le plus précis est le plus informatif.\end{attention}
\clearpage
\phantomsection\label{cours-16-3}
\essSection{3}{Pentagones et hexagones}
\begin{definition}Un pentagone a 5 côtés et 5 sommets ; un hexagone, 6 côtés et 6 sommets. Un polygone est régulier quand tous ses côtés et tous ses angles sont égaux : l’hexagone des alvéoles d’abeilles est régulier.\end{definition}
\label{fig-cours-16-3-cours-hexagone}\figureGeo{cours-hexagone}

\begin{savoirfaire}[L’hexagone au compas]Construire un hexagone régulier : tracer un cercle ; garder la même ouverture de compas (le rayon) et la reporter six fois sur le cercle ; joindre les six points. Le tour de cercle se ferme exactement.\end{savoirfaire}
\begin{exemple}Le côté de cet hexagone régulier est égal au rayon du cercle : un hexagone construit dans un cercle de rayon 4 cm a des côtés de 4 cm et un périmètre de 24 cm.\end{exemple}
\clearpage
\phantomsection\label{cours-16-4}
\essSection{4}{Construire et justifier}
\begin{savoirfaire}[Les instruments tranchent]Avant de nommer une figure, vérifier ses propriétés aux instruments : équerre pour les angles droits, compas ou règle graduée pour les longueurs, règle et équerre pour le parallélisme. L’allure ne suffit jamais.\end{savoirfaire}
\begin{exemple}Un trapèze rectangle de bases 6 cm et 4 cm et de hauteur 3 cm se construit comme un rectangle auquel on incline un côté : tracer la base, deux perpendiculaires, reporter 6 cm et 4 cm, joindre.\end{exemple}
\begin{propriete}Une figure composée se décompose en figures simples : un trapèze rectangle se découpe en un rectangle et un triangle rectangle. Ce découpage sert pour les contours, les aires et les programmes de construction.\end{propriete}
\begin{attention}Les diagonales du rectangle ont la même longueur et se coupent en leur milieu : cette observation, vérifiée à la règle, contrôle efficacement une construction.\end{attention}

\clearpage\rubrique{Découverte}
\begin{blocLivre}{ex-16-1}\begin{exercice}[1]{Compter les côtés}
Nommer chacun des quatre polygones dessinés : triangle, quadrilatère, pentagone ou hexagone.
\label{fig-ex-16-1}\figureGeo{tri-1}
\end{exercice}\end{blocLivre}
\begin{blocLivre}{ex-16-2}\begin{exercice}[1]{Classer les triangles}
Classer les quatre triangles dessinés : rectangle, isocèle, équilatéral ou quelconque, en lisant les codes.
\label{fig-ex-16-2}\figureGeo{tri-2}
\end{exercice}\end{blocLivre}
\begin{blocLivre}{ex-16-3}\begin{exercice}[1]{L’isocèle au compas}
Construire un triangle isocèle de base 4 cm et de côtés égaux de 6 cm. Garder les arcs.
\label{fig-ex-16-3}\figureGeo{tri-3}
\end{exercice}\end{blocLivre}
\begin{blocLivre}{ex-16-4}\begin{exercice}[1]{Trapèze ou pas ?}
Parmi les quatre quadrilatères dessinés, lesquels sont des trapèzes ? Justifier avec le parallélisme.
\label{fig-ex-16-4}\figureGeo{tri-4}
\end{exercice}\end{blocLivre}
\begin{blocLivre}{ex-16-5}\begin{exercice}[1]{La carte d’identité}
Pour le trapèze rectangle dessiné : combien d’angles droits ? Quels côtés sont parallèles ?
\label{fig-ex-16-5}\figureGeo{tri-5}
\end{exercice}\end{blocLivre}

\rubrique{Entraînement}
\begin{blocLivre}{ex-16-6}\begin{exercice}[2]{Le rectangle contrôlé}
Construire un rectangle de 7 cm sur 4 cm, tracer ses diagonales et les mesurer. Que constate-t-on ?
\label{fig-ex-16-6}\figureGeo{tri-6}
\end{exercice}\end{blocLivre}
\begin{blocLivre}{ex-16-7}\begin{exercice}[2]{Le carré incliné}
Construire un carré de côté 5 cm dont un côté est déjà tracé en biais. L’orientation change-t-elle la nature de la figure ?
\label{fig-ex-16-7}\figureGeo{tri-7}
\end{exercice}\end{blocLivre}
\begin{blocLivre}{ex-16-8}\begin{exercice}[2]{Le losange au compas}
Construire un losange de côté 4 cm à partir de l’angle donné au croquis, au compas seulement.
\label{fig-ex-16-8}\figureGeo{tri-8}
\end{exercice}\end{blocLivre}
\begin{blocLivre}{ex-16-9}\begin{exercice}[2]{L’hexagone régulier}
Construire un hexagone régulier dans un cercle de rayon 4 cm, en reportant le rayon au compas. Donner son périmètre.
\label{fig-ex-16-9}\figureGeo{tri-9}
\end{exercice}\end{blocLivre}
\begin{blocLivre}{ex-16-10}\begin{exercice}[2]{Régulier ou non ?}
Deux pentagones sont dessinés : lequel est régulier ? Justifier.
\label{fig-ex-16-10}\figureGeo{tri-10}
\end{exercice}\end{blocLivre}

\rubrique{Problèmes}
\begin{blocLivre}{ex-16-11}\begin{exercice}[3]{Le trapèze décomposé}
Montrer que le trapèze rectangle dessiné (bases 6 cm et 4 cm, hauteur 3 cm) se découpe en un rectangle et un triangle rectangle. Donner les dimensions de chacun.
\label{fig-ex-16-11}\figureGeo{tri-11}
\end{exercice}\end{blocLivre}
\begin{blocLivre}{ex-16-12}\begin{exercice}[3]{Le programme de l’hexagone}
Écrire un programme de construction de l’hexagone régulier de côté 3 cm.
\label{fig-ex-16-12}\figureGeo{tri-12}
\end{exercice}\end{blocLivre}
\begin{blocLivre}{ex-16-13}\begin{exercice}[3]{La table basse}
Le plateau d’une table a la forme d’un trapèze : bases de 90 cm et 60 cm, deux autres côtés de 40 cm chacun. Calculer la longueur de la baguette de chant qui fait tout le tour.
\label{fig-ex-16-13}\figureGeo{tri-13}
\end{exercice}\end{blocLivre}
\begin{blocLivre}{ex-16-14}\begin{exercice}[3]{Les alvéoles}
Les abeilles pavent leurs rayons d’hexagones réguliers. Combien de côtés chaque alvéole partage-t-elle avec ses voisines quand elle est entourée de toutes parts ? Quel est l’intérêt de ce partage ?
\label{fig-ex-16-14}\figureGeo{tri-14}
\end{exercice}\end{blocLivre}
\begin{blocLivre}{ex-16-15}\begin{exercice}[3]{Qui est qui ?}
Vrai ou faux : « Tout carré est un losange » ; « Tout losange est un carré » ; « Tout rectangle a ses côtés opposés parallèles ». Justifier chaque réponse.
\label{fig-ex-16-15}\figureGeo{tri-15}
\end{exercice}\end{blocLivre}

\clearpage
